\section{Decis\~oes de projeto (D1 a D17)}\label{sec:decisoes}

O plano registra 17 decis\~oes. As dez primeiras foram tomadas em 03/11/2026, ao validar o plano.
As demais surgiram durante a implementa\c{c}\~ao, quando um teste ou experimento mostrou que a
ideia inicial n\~ao funcionava. Entender o \emph{porqu\^e} de cada uma \'e o que permite defender
o trabalho.

{\small
\begin{longtable}{@{}p{0.7cm}p{2.5cm}p{5.6cm}p{6.6cm}@{}}
\toprule
\# & Tema & Decis\~ao & Por qu\^e \\ \midrule \endhead
D1 & Painel & ZTROON ZTP-160M; o KC200GT fica s\'o como benchmark da extra\c{c}\~ao
  & Valida o m\'etodo contra um resultado publicado (Villalva, 2009) e aplica ao painel real \\
D2 & Ferramentas & MATLAB base; depois, tamb\'em a Optimization Toolbox s\'o para o \texttt{fsolve} & Inicialmente sem toolboxes (LM pr\'oprio). Revisto em 10/11/2026: o \texttt{fsolve} em $\log_{10}$ de $R_s$ e $R_p$ d\'a os mesmos par\^ametros \\
D3 & Bateria & 12\,V constante (fonte ideal)
  & Sem sombra $V_{mp}\approx19{,}6$\,V $>12$\,V: o conversor abaixa. Com sombra o pico global cai para 7 a 13\,V: precisa elevar. O buck-boost \'e necess\'ario justamente para o MPPT global \\
D4 & Conversor & $f_s=20$\,kHz, $L\approx57\,\mu$H, $C\approx1{,}3$\,mF, perdas de cat\'alogo
  & 20\,kHz \'e inaud\'ivel e tem perdas de chaveamento pequenas nessa pot\^encia (o trabalho anterior usou 2\,kHz, com indutor 10 vezes maior) \\
D5 & EN 50530 & Tabelas B.2 a B.4 transcritas de um relat\'orio de ensaio T\"UV Rheinland
  & O texto da norma \'e pago; os valores conferem com as dura\c{c}\~oes impressas \\
D6 & Dia real & PVGIS hor\'ario de Florian\'opolis, com nuvens e sombra de poste sint\'eticas e declaradas
  & Dados hor\'arios n\~ao t\^em transit\'orios r\'apidos \\
D7 & Relat\'orio & \LaTeX, estilo do trabalho anterior & Continuidade com a disciplina \\
D8 & Valida\c{c}\~ao do painel & Sem curvas I-V (o datasheet n\~ao tem): STC, coeficiente t\'ermico de $P_{max}$, NMOT e KC200GT
  & Usa s\'o informa\c{c}\~ao publicada e declara a limita\c{c}\~ao \\
D9 & 2 diodos & Conven\c{c}\~ao de Ishaque: $a_1=1$, $a_2=1{,}2$, $I_{01}=I_{02}$; lei t\'ermica efetiva com $E_g$ ajustado pelo $K_v$
  & Com $E_g=1{,}12$\,eV o modelo d\'a $K_v=-0{,}106$\,V/K contra $-0{,}123$ medido; e o datasheet n\~ao identifica $I_{02}$ \\
D10 & $R_p$ & $R_p(G)=R_{p,ref}\,G_n/G$ (De Soto)
  & Com $R_p$ constante o ajuste previa perda de 26\,\% em 200\,W/m$^2$, contra 3 a 5\,\% t\'ipicos \\
D11 & T\'ermico & Convec\c{c}\~ao de Watmuff, radia\c{c}\~ao frente$\to$c\'eu e verso$\to$solo, capacidade pela massa do datasheet
  & McAdams j\'a inclui radia\c{c}\~ao e a contava duas vezes. Com as corre\c{c}\~oes: $k_h=0{,}79$ e $\tau\approx7$\,min \\
D12 & Projeto do PI & Varrer $f_c$ e a posi\c{c}\~ao do zero $r=\omega_z/\omega_c$, com MF$\ge45^\circ$, MG$\ge6$\,dB; escolher a maior $f_c$
  & A regra inicial (zero em $\omega_c/5$) falha: a resson\^ancia LC ($\sim$240\,Hz, $\zeta\approx0{,}16$) cruza 0\,dB com fase $\approx-150^\circ$ \\
D13 & Robustez e $\Ta$ & Margens em todos os pontos $\times$ cantos de $L$ e $C$ ($\pm20$\,\%) e ponto quente; $\Ta=1{,}5\,t_s$ no pior caso
  & Com $L+20$\,\% a MG caiu a 4,6\,dB; o MPPT tem de esperar o transit\'orio do PI \\
D14 & Par\^ametros do MPPT & InC: $\varepsilon=\tfrac12|ds/dV|\delta$. PSO: part\'iculas em $k\,V_{mp}/N_{sub}$ e $0{,}9V_{oc}$
  & $\varepsilon=0{,}002$\,S nunca era atingido (InC oscilava como P\&O); com part\'iculas espa\c{c}adas linearmente o PSO errava o pico em $\sim$1/3 das sementes \\
D15 & Din\^amica & EDO reduzida de 4 estados, com a corrente do painel resolvida por Newton a cada avalia\c{c}\~ao
  & A DAE com matriz de massa falhou em 9 de 36 casos; a EDO: 0 falhas, mesma efici\^encia em 4 d\'igitos, 2 a 4 vezes mais lenta \\
D16 & Quase-est\'atico & Aceito com erro $\le0{,}4$\,\% (P\&O, InC) e $-0{,}9$ a $-1{,}1$\,\% (PSO); or\'aculo com $dt=1$\,s
  & Com $dt=\Ta$ o or\'aculo custa $\sim$2500\,s de CPU por hora simulada; com 1\,s, $\sim$65\,s \\
D17 & Limita\c{c}\~oes & Documentar como resultado, sem alterar os algoritmos
  & Os tr\^es efeitos s\'o aparecem com o modelo t\'ermico e o dia real; corrigir exigiria rodar E6 e E7 de novo (horas) \\
\bottomrule
\end{longtable}}

\begin{resumo}
D1 a D8 s\~ao escolhas de escopo. D9 a D17 s\~ao \emph{corre\c{c}\~oes de rumo}: o plano dizia uma
coisa, o teste mostrou outra, e a decis\~ao registra o que foi trocado e por qu\^e. O padr\~ao se
repete: primeiro a hip\'otese mais simples, depois a medi\c{c}\~ao que a derruba, depois a
corre\c{c}\~ao documentada.
\end{resumo}
